\section{Application templates and repaired difficulty interfaces}
\label{sec:applications}

The framework is intended to organize many branches of mathematics without
pretending that their deep quantitative content is automatic.  This section
shows how the common compiler layer applies and identifies exactly what remains
theorem-specific.

\subsection{Finite algebra and linear algebra}

For a fixed matrix dimension, every ring expression, determinant, minor,
characteristic polynomial, and matrix product is a finite PR computation.
The PCOG decomposition is:

\begin{enumerate}
\item exact nodes for coefficient and matrix arithmetic;
\item bounded folds for finite sums and products;
\item identity-certificate nodes for polynomial equalities;
\item trace nodes for elimination or iterative reduction;
\item observation nodes for invariants stable under basis or witness choices.
\end{enumerate}

This architecture covers constructive versions of Cayley--Hamilton, finite
module computations, rational canonical data, determinant identities, and
bounded Gröbner-style reductions.  The cleanest Cayley--Hamilton route uses a
finite polynomial identity over the generic matrix and then substitutes the
input entries.  The proof obligation is coefficientwise and bounded by the
fixed dimension.  No eigenvalues are needed.

For variable dimension \(n\), the algorithms remain PR because table sizes,
permutation enumeration, and recursion depths are PR in \(n\).  The object
proof uses one uniform invariant over coded matrices.  The challenge is not
computability but the size and organization of algebraic lemmas.  Proof-
carrying nodes permit each identity to be checked independently.

\subsection{Polynomial systems and algebraic geometry}

A finite affine polynomial system is coded by coefficient tables and degree
bounds.  PCOG supports the following layers.

\paragraph{Exact local algebra.}
Evaluation, derivatives, Jacobians, resultants at fixed degree, Sylvester
matrices, finite quotient-ring multiplication tables, and coefficientwise
identities are exact nodes.  A certificate that a finite span contains one and
is closed under multiplication by coordinates yields a finite-dimensional
algebra bound by straightforward matrix checks.

\paragraph{Filtered lifting.}
Given a seed and a two-sided inverse Jacobian modulo the first filtration
level, a fixed number of Newton-style updates in a nilpotent filtered ring is a
trace node.  Residual improvement is the invariant.  The coherent-methods
packet gives a broad human theorem for this situation.

\paragraph{Root and branch observations.}
Simple roots with a derivative apartness witness produce regular names by a
contractive or Newton trace.  Singular fibers produce finite candidate clouds.
Exact branch labels are exported only after a separation theorem.

\paragraph{Finite covers and sections.}
A cover is a finite relational output.  Restriction maps and overlap checks
are bounded.  A global construction can be represented by compatible local
observations and backward viability.  This gives a finitary interpretation of
many sheaf-like arguments without quantifying over an actual infinite sheaf.

\paragraph{Projective and graded data.}
At each degree bound \(d\), homogeneous ideals, syzygies, Hilbert-function
values, and finite graded pieces are coded tables.  A theorem requiring all
degrees must supply a PR regularity or stabilization bound, or remain a family
of finite statements indexed by \(d\).  The framework makes the missing
regularity producer explicit.

This is a plausible route to substantial algebraic geometry in iPRA: focus on
finite presentations, explicit degree bounds, local certificates, and
quantitative stabilization.  It does not collapse abstract existence theorems
without effective bounds into finite computation.

\subsection{Analysis and differential equations}

Constructive analysis in iPRA is naturally expressed through regular rational
names and finite error certificates.

\paragraph{Series and products.}
A partial-sum node is a trace whose invariant records the exact recurrence and
an error bound.  A tail estimate node turns a requested precision into a
truncation demand.  The finalizer combines the partial sum and tail bound.

\paragraph{Differential equations.}
For rational polynomial ODEs on a certified bounded region, Euler or Taylor
schemes are trace nodes.  The invariant records state bounds, residual bounds,
and error recurrence.  Conservation laws are exact polynomial observation
nodes.  A local solution theorem needs explicit Lipschitz and region-invariance
certificates; larger horizons need a separate continuation resource.

\paragraph{Integration.}
Riemann sums, quadrature rules, and interval enclosures are finite tables.  A
modulus of continuity or derivative bound determines mesh demand.  The
integral is a stable observable.  Improper integrals require a tail modulus.

\paragraph{Special functions.}
Instead of importing a completed transcendental object, define a PR rational
name by a recurrence or ODE trace and prove its contracts.  Identities between
special functions become comparison nodes or polynomial conservation laws.

A recurrent danger is to prove a convergence theorem in ambient Lean and then
use its limit as if it were an iPRA term.  PCOG prevents this: the exported
object is always the PR name plus first-order error theorem.

\subsection{Topology and finite geometry}

General point-set topology is not a primitive-recursive data type, but many
usable topological constructions have finite quantitative presentations.

\paragraph{Finite spaces and complexes.}
Finite preorders, simplicial complexes, chain groups, boundary matrices, and
homology up to a fixed degree are exact finite data.  Reduction algorithms are
bounded traces.  Invariants such as Betti numbers are stable observations.

\paragraph{Compact metric presentations.}
A compact object can be represented by PR finite nets with coverage proofs and
refinement maps.  Uniformly continuous maps carry explicit moduli.  Products,
images, finite covers, and approximate fixed-point searches become demand
graphs.  Exact points require additional uniqueness or contraction data.

\paragraph{Homotopy data.}
At a fixed mesh, homotopies are finite tables with local compatibility
checks.  Passing to arbitrary precision needs a refinement theorem and a PR
modulus.  A theorem that only asserts abstract compactness without such data
cannot be compiled directly.

The framework thus interprets topology through finite observations rather than
through arbitrary collections of open sets.  This is not merely a coding
restriction; it identifies which topological theorems have quantitative
computational content at iPRA strength.

\subsection{Category-theoretic organization}

Category theory is useful as an external language for composing finite
presentations.  Fix a class of presentations and proof-carrying morphisms.
Objects are code/validity/restriction packages; a morphism is a PR transformer
with demand and local soundness proof.  Identity and composition are given by
identity functions and demand composition.  Associativity holds by equality
of the emitted PR syntax up to provable extensional equality.

At the internal first-order level, one need not formalize a universe of all
categories.  Each theorem uses a fixed finite diagram.  Limits, colimits,
adjunction components, or Yoneda-style evaluations become finite graph
constructions for that diagram.  The categorical layer organizes the host
compiler, while the emitted theorem remains ordinary iPRA syntax.

This distinction avoids a common overreach: a categorical proof can guide the
compiler without requiring iPRA to quantify over arbitrary functors or natural
transformations as higher-order objects.

\subsection{Number theory and bounded arithmetic}

Many elementary and computational number-theoretic arguments fit directly.

\begin{itemize}
\item Euclidean algorithms and continued fractions use ranked trace nodes.
\item Primality and factorization certificates use bounded verification and
finite search; a total factorization producer is PR by bounded trial division,
though inefficient.
\item Congruence lifting and Chinese remainder calculations are exact nodes.
\item Finite sums over divisors, residue classes, or combinatorial partitions
are bounded folds.
\item Diophantine searches are complete once a height bound or finite cover is
supplied.
\item Explicit analytic number theory can use rational interval arithmetic,
finite prime tables, and proved remainder estimates.
\end{itemize}

The hard part of advanced prime-distribution theorems is not representing
finite sums.  It is supplying uniform effective estimates with a strength and
logical form available in iPRA.  PCOG can verify and compose such estimates;
it cannot infer them from their classical existence.

\subsection{Repairing the Siegel and Faltings interfaces}

The proof collection contains zero-count, norm, section, and finite-cover
arguments relevant to Siegel and Faltings.  These solve genuine local
obligations but not the general effective theorem.  The PCOG framework gives a
precise conditional endpoint.

Let \(C\) be a coded curve or variety and \(P(C,x)\) the decidable predicate
that \(x\) is a rational or integral point in normalized code.  Suppose a node
produces an explicit height bound \(B(C)\) with an embedded proof
\[
 P(C,x)\longrightarrow \height(x)\leq B(C).
\]
Then Theorem~\ref{thm:height-to-list} produces a complete finite point list.
All algebraic verification, duplicate removal, and completeness bookkeeping
are routine PCOG nodes.

Alternatively, suppose a node produces a finite auxiliary cover and every
point lies in some section for which a complete bounded list is available.
Finite auxiliary exhaustion produces the point list without selecting a
section per point.

This resolves the \emph{post-bound} formalization problem completely.  It does
not provide the missing general height or cover producer.  The report therefore
recommends that future Siegel/Faltings packets state their unproved interface
as exactly one of:

\begin{enumerate}
\item a PR height-bound symbol plus its embedded theorem;
\item a PR complete-list symbol plus soundness and completeness;
\item a finite cover producer plus coverage and section-completeness proofs.
\end{enumerate}

A cardinality theorem alone should not be presented as enough for enumeration.

\subsection{Repairing the Mordell--Weil interface}

A constructive Mordell--Weil theorem can be decomposed into:

\begin{enumerate}
\item exact coded group-law nodes on rational points;
\item a finite descent map and local solvability certificates;
\item a finite covering list of relevant cosets or homogeneous spaces;
\item a height-reduction theorem and bounded descent fuel;
\item a finite list of generators and an embedded proof that every point
reduces to their span.
\end{enumerate}

PCOG handles steps 1, 2, 4, and the verification side of 5.  The unresolved
general issue is a uniform effective producer for the covering list and the
necessary arithmetic information.  If that list is supplied, backward
viability can retain all locally compatible lifts and finite auxiliary
exhaustion can avoid arbitrary selectors.

The public theorem should distinguish:
\begin{itemize}
\item a conditional generator theorem from a supplied covering package;
\item a complete theorem for explicitly bounded families;
\item the still-missing uniform general producer.
\end{itemize}

\subsection{Repairing singular branches and exact orientation choices}

For a simple analytic branch, a root seed and derivative apartness witness
support a Newton or coefficient recurrence.  For a singular point, arbitrary
Cauchy coefficient names do not uniformly determine exact Puiseux branch
coefficients or labels.  The relational repair is:

\begin{enumerate}
\item at each truncation order, compute a finite cloud of candidate branch
jets satisfying a residual checker;
\item connect compatible truncations by an extension relation;
\item use finite-horizon viability to remove dead ends;
\item export symmetric or set-valued observations through collisions;
\item collapse to one branch only with an isolation certificate.
\end{enumerate}

The same pattern applies to choosing an exact rotation axis at a symmetric or
zero input.  A total exact side/axis selector can imply LLPO-like principles.
A projective line, unoriented axis, finite candidate set, or gap-guarded exact
choice is the correct output.

\subsection{Prime-distribution packets}

Dirichlet, Chebotarev, Brun--Titchmarsh, Siegel--Walfisz, and
Bombieri--Vinogradov involve several layers:

\begin{itemize}
\item finite arithmetic functions and sums, which are straightforward PR
nodes;
\item rational approximations to analytic quantities, which need moduli and
error traces;
\item zero-free regions, density estimates, or exceptional-character bounds,
which are substantive quantitative producers;
\item final finite inequalities and bounded searches, which PCOG can compile.
\end{itemize}

A responsible iPRA formalization should present a dependency graph whose
unproved nodes are the exact estimates, with all constants and ranges explicit.
The framework then ensures that once those nodes are proved, no further hidden
compactness, choice, or limit principle remains in the assembly.

\subsection{A difficulty matrix}

\begin{longtable}{@{}p{0.19\textwidth}p{0.28\textwidth}p{0.25\textwidth}p{0.20\textwidth}@{}}
\toprule
Area & Finite interface & PCOG resolution & Remaining producer \\
\midrule
\endhead
Simple analytic root & seed, derivative gap, coefficient/interval checker &
trace plus guarded inverse; regular name & explicit local radius and bounds \\
Singular branch & finite jet clouds and extension relation & viability,
shrinking observations, gap collapse & candidate completeness and separation \\
Matrix spectrum & rational diagonalization certificates & independent-witness
comparison, ordered eigenvalue observation & matrix-specific residual estimates \\
Finite topology & nets, covers, chain tables & demand composition, bounded
homology traces & quantitative refinement/coverage \\
Siegel/Faltings & point predicate, height/list/cover interface & complete
bounded enumeration and verification & general effective height or cover \\
Mordell--Weil & finite descent/cover package & relational lifts, bounded descent,
generator verification & uniform covering/generator producer \\
Prime distribution & finite sums plus explicit estimates & graph composition,
finite windows, rational inequalities & deep zero-density/zero-free estimates \\
Metastability & bounded stable-window predicate & finite union/search node & a
separate principle if a global modulus is desired \\
Exact branch/sign & finite candidate relation & exact collapse only with gap &
apartness or separation certificate \\
\bottomrule
\end{longtable}

The matrix clarifies what ``a more general framework'' can honestly achieve.
It can make the finite proof pipeline uniform and complete.  It cannot replace
subject-specific quantitative theorems whose absence is the mathematical
difficulty.
